\newpage
\phantomsection
\label{prf:absolute-value-preserves-continuity}

\begin{remark*}[Return]
\hyperref[thm:absolute-value-preserves-continuity]{Return to Theorem}
\end{remark*}

\begin{theorem*}[Absolute Value Preserves Continuity]
Let $A \subseteq \mathbb{R}$, let $f : A \to \mathbb{R}$, and let $|f|$ be defined by $|f|(x):=|f(x)|$ for $x\in A$. If $f$ is continuous at $c\in A$ then $|f|$ is continuous at $c$; if $f$ is continuous on $A$ then $|f|$ is continuous on $A$.
\end{theorem*}

\begin{proof}
\textbf{Professional Standard Proof.}~\\
Let $A \subseteq \mathbb{R}$, let $f : A \to \mathbb{R}$, and let $|f|(x):=|f(x)|$ for $x\in A$.

First, suppose $f$ is continuous at $c\in A$. Let $\varepsilon>0$. Since $f$ is continuous at $c$, there exists $\delta>0$ such that $|x-c|<\delta$ implies $|f(x)-f(c)|<\varepsilon$ for all $x\in A$. For any $x\in A$ with $|x-c|<\delta$, the reverse triangle inequality gives $\bigl||f(x)|-|f(c)|\bigr| \le |f(x)-f(c)|<\varepsilon$. Since $|f|(x)=|f(x)|$ and $|f|(c)=|f(c)|$, this shows $\bigl||f|(x)-|f|(c)\bigr|<\varepsilon$. Therefore $|f|$ is continuous at $c$.

Second, if $f$ is continuous on $A$, then $f$ is continuous at every point $c\in A$. By the first part, $|f|$ is continuous at every point $c\in A$. Therefore $|f|$ is continuous on $A$.
\end{proof}

\begin{proof}
\textbf{Detailed Learning Proof.}~\\

\textbf{Step 1.} Establish continuity at a point. Assume $f$ is continuous at $c\in A$ and let $\varepsilon>0$ be given. By the definition of continuity, there exists $\delta>0$ such that for every $x\in A$, if $|x-c|<\delta$ then $|f(x)-f(c)|<\varepsilon$. This gives the foundation for proving that $|f|$ is continuous at $c$.

\textbf{Step 2.} Apply the reverse triangle inequality. For any $x\in A$ satisfying $|x-c|<\delta$, the reverse triangle inequality states that $\bigl||f(x)|-|f(c)|\bigr| \le |f(x)-f(c)|$. Combined with the continuity condition from Step 1, this yields $\bigl||f(x)|-|f(c)|\bigr| < \varepsilon$.

\textbf{Step 3.} Translate to the absolute value function. Since $|f|(x)=|f(x)|$ and $|f|(c)=|f(c)|$ by definition, the inequality from Step 2 becomes $\bigl||f|(x)-|f|(c)\bigr|<\varepsilon$ whenever $x\in A$ and $|x-c|<\delta$. This establishes that $|f|$ is continuous at $c$.

\textbf{Step 4.} Extend to continuity on the domain. If $f$ is continuous on $A$, then $f$ is continuous at every point $c\in A$. By Steps 1-3, $|f|$ is continuous at every such point $c$. Since this holds for arbitrary $c\in A$, the function $|f|$ is continuous on $A$.
\end{proof}

\begin{remark*}[Proof structure]
This proof employs a direct approach using the reverse triangle inequality as the key tool. The reverse triangle inequality $\bigl||a|-|b|\bigr| \le |a-b|$ transforms the problem of controlling differences of absolute values into controlling differences of the original function values, which is exactly what continuity provides.
\end{remark*}

\begin{remark*}[Dependencies]~\\
\begin{itemize}
  \item \hyperref[def:continuous-at-point]{Definition of Continuity at a Point}
  \item \hyperref[def:continuity-on-set]{Definition of Continuity on a Set}
  \item \hyperref[thm:reverse-triangle-inequality]{Reverse Triangle Inequality}
\end{itemize}
\end{remark*}
